%%%
\section{Python apps tests}
%%%
\subsection{Overview}
In addition to the SW2D executables, coded in C++ and compiled, Python scripts are available in utils/python/apps. \\
These scripts are auxiliary to the software and are therefore not subject to mr-tests required for branch merging.
However, there are test cases intended to test their proper functioning, which can be used when necessary (for example, when making changes to these scripts).

\subsection{Functioning}

For each app, a single Python script runs all its tests and displays the results in the console.
Subsection \ref{py_tests_usage} explains how to run these tests.

\subsubsection{\swConverterReverse}
These tests run the \swConverterReverse for each case and then compare the outputs with reference.\\

Since the \swConverterReverse is supposed to reproduce the input files used to create a \file{geometry.sw2d} file, the reference files for a given test case are the ones used to create the \file{geometry.sw2d} for that case.
The tests verify that the \swConverterReverse runs without error and the outputs are exactly identical to the reference.\\

At the end of the tests, a table is displayed in the console.
It contains one row per test case and one column per element tested (\swConverterReverse successful execution, correct results).
All fields must be \texttt{Success}.

\subsubsection{\swNSPreProcess}
These tests run the \swNSPreProcess for each case and then compare the outputs with reference.\\

Since the \swNSPreProcess reads nodetrings in a .2dm file and create the .bc and qsections.txt file, the reference files are hand-made by developers.
The tests verify that the \swNSPreProcess runs without error and the outputs are exactly identical to the reference. \\

At the end of the tests, a table is displayed in the console.
It contains one row per test case and one column per element tested (\swNSPreProcess successful execution, correct results).
All fields must be “Success.”

\subsubsection{\swOutputConverter}
These tests run the \swConverter, the \swSolver and the \swOutputConverter for each case and then compare the outputs with reference.\\

The \swOutputConverter can convert a binary file into a netcdf file or text files. In these tests, the file is converted into text, and the reference files for a test case are the solver text outputs for that test.
The tests verify that \swConverter, \swSolver and \swOutputConverter run without error and the outputs are identical to the reference with a tolerance defined in the test script.
This margin of error is present because converting doubles to binary and then to text can lead to minor approximations.\\

At the end of the tests, a table is displayed in the console.
It contains one row per test case and one column per element tested (\swConverter, \swSolver, \swOutputConverter successful execution, correct results) and a column for computation time.
All fields must be “Success.”

\subsection{Usage} \label{py_tests_usage}
To run all the test cases dedicated to a python app, follow these steps :
\begin{itemize}
    \item Open a shell/prompt.
    \item Navigate to \file{utils/python/apps/tests/\textlangle app\_name\textrangle} where \textlangle app\_name\textrangle  is the name of the app you want to test.
    \item Type :
    \begin{verbatim}
	python <app_name>_tests.py
	\end{verbatim}
    This will run the tests and display the result table in the console.
\end{itemize}